\source{Лекции, с. 21–27; Карпов, §9.2, с. 425–428; детали Бине–Коши восстановлены составителем}
\begin{definition}[Матрица Кирхгофа]
Для неориентированного графа без петель, в котором допускаются кратные рёбра,
лапласиан (матрица Кирхгофа) равен \(K=D-A\), где \(D=\diag(d(v_1),\ldots,d(v_n))\).
Иными словами, \(K_{ii}=d(v_i)\), \(K_{ij}=-a_{ij}\) при \(i\ne j\).
Через \(K^{(r)}\) обозначим матрицу после удаления строки и столбца \(r\),
через \(\tau(G)\) — число остовных деревьев.
\end{definition}
\begin{definition}[Матрицы инцидентности]
Неориентированная матрица инцидентности имеет строки-вершины и столбцы-рёбра:
\(I_{ve}=1\), если \(v\) — конец \(e\), иначе \(0\).
Ориентированная матрица меняет у конца дуги единицу на \(-1\),
а у начала оставляет \(+1\). В доказательстве обозначаем её \(B\),
чтобы не путать с единичной матрицей.
\end{definition}
\begin{lemma}[Матрица инцидентности]
Произвольно ориентируем рёбра и составим матрицу \(B\) размера \(n\times m\):
в столбце дуги \(u\to v\) стоят \(+1\) у \(u\), \(-1\) у \(v\), прочие элементы нулевые.
Тогда \(BB^{\mathsf T}=K\). Если удалить строку \(r\), полученная матрица \(B_r\)
удовлетворяет \(B_rB_r^{\mathsf T}=K^{(r)}\).
\end{lemma}
\begin{proof}
Скалярный квадрат строки равен числу инцидентных рёбер — степени вершины.
Скалярное произведение разных строк получает вклад \(-1\) ровно от рёбер
между соответствующими вершинами. Это и есть элементы \(K\).
Удаление строки в обоих множителях даёт нужную главную подматрицу.
\end{proof}
\begin{lemma}[Миноры инцидентности]
Пусть \(S\) состоит из \(n-1\) рёбер и \(B_{r,S}\) — соответствующая
квадратная подматрица \(B_r\). Тогда
\[
 \det B_{r,S}=\begin{cases}\pm1,&(V,S)\text{ — дерево},\\0,&\text{иначе}.\end{cases}
\]
\end{lemma}
\begin{proof}
Если \((V,S)\) не дерево, при \(|S|=n-1\) он несвязен по билету 6.
Выберем компоненту, не содержащую \(r\). Сумма строк её вершин
в \(B_{r,S}\) равна нулю: у каждого выбранного ребра оба конца в этой
компоненте либо оба вне её. Строки линейно зависимы, определитель нулевой.
Для дерева доказываем модуль определителя индукцией: при \(n=1\)
определитель пустой матрицы равен 1. При \(n>1\) есть лист \(x\ne r\).
В его строке единственный ненулевой элемент \(\pm1\). Разложение по этой
строке сводит определитель к меньшему дереву с удалённым листом.
\end{proof}
\begin{lemma}[Бине–Коши в нужном виде]
Для матрицы \(M\) размера \(p\times q\), \(p\le q\),
\[
 \det(MM^{\mathsf T})=\sum_{S\subseteq[1,q],\ |S|=p}(\det M_{[:,S]})^2.
\]
\end{lemma}
\begin{proof}
Более общо, для \(X\) размера \(p\times q\) и \(Y\) размера \(q\times p\)
разложим каждый столбец \(XY\) как линейную комбинацию столбцов \(X\).
Мультилинейность определителя даёт сумму по \(p\)-кам индексов.
Повторяющиеся индексы дают нуль из-за равных столбцов. Для фиксированного
набора различных индексов \(S\) сумма по всем их перестановкам равна
\(\det X_{[:,S]}\det Y_{[S,:]}\): знаки перестановок возникают при
перестановке столбцов \(X\), а сумма произведений элементов \(Y\)
со знаками есть её определитель. При \(Y=X^{\mathsf T}\) оба определителя равны.
\end{proof}
\begin{theorem}[Кирхгоф]
Для любого \(r\) и \(n\ge1\)
\[
 \boxed{\tau(G)=\det K^{(r)}}.
\]
\end{theorem}
\begin{proof}
Если \(m<n-1\), остовных деревьев нет; ранг \(B_rB_r^{\mathsf T}\)
не больше \(m<n-1\), поэтому определитель нулевой.
Иначе по трём леммам
\[
 \det K^{(r)}=\det(B_rB_r^{\mathsf T})
 =\sum_{|S|=n-1}(\det B_{r,S})^2=\sum_{(V,S)\text{ — дерево}}1.
\]
Для \(n=1\) обе стороны равны 1. Доказательство также охватывает несвязные графы:
каждый минор в сумме нулевой и \(\tau(G)=0\).
\end{proof}
\begin{cor}[Все алгебраические дополнения]
При \(n\ge2\) для любых \(i,j\)
\( (-1)^{i+j}\det K_{\widehat i,\widehat j}=\tau(G)\).
\end{cor}
\begin{proof}
Если \(G\) связен, основной минор положителен, а \(K\boldsymbol1=0\).
Поэтому \(\rank K=n-1\) и ядро состоит из постоянных векторов.
Из \(K\operatorname{adj}K=0\) каждый столбец присоединённой матрицы
постоянен; её симметрия делает все элементы одинаковыми.
На диагонали они равны \(\tau(G)\). Если компонент не менее двух,
их индикаторы — независимые векторы ядра, так что \(\rank K\le n-2\)
и все миноры порядка \(n-1\) нулевые.
\end{proof}
\begin{example}[Лекционный граф на четырёх вершинах]
Возьмём рёбра \(12,23,34,41,13\): квадрат с одной диагональю.
Тогда
\[
K=\begin{pmatrix}3&-1&-1&-1\\-1&2&-1&0\\-1&-1&3&-1\\-1&0&-1&2\end{pmatrix},
\quad \tau(G)=\det\begin{pmatrix}2&-1&0\\-1&3&-1\\0&-1&2\end{pmatrix}=8.
\]
В предварительном вычислении на с. 21 рукописи у матрицы стоят обратные
знаки; здесь используется последующее определение \(K=D-A\), поэтому
определитель главного минора положителен.
\end{example}
\begin{cor}[Удаление–стягивание для числа остовов]
Для непетлевого ребра \(e\)
\[
\tau(G)=\tau(G-e)+\tau(G/e).
\]
\end{cor}
\begin{proof}
Остовы, не содержащие \(e\), — в точности остовы \(G-e\).
Стягивание \(e\) даёт биекцию между остовами, содержащими \(e\),
и остовами \(G/e\): обратная операция раскрывает стянутую вершину
и возвращает \(e\). Связность и отсутствие циклов при этом сохраняются.
После стягивания петли удаляются, но кратные рёбра сохраняются различимыми.
\end{proof}
